\documentclass[12pt]{article}
\usepackage[utf8]{inputenc}
\usepackage[T1]{fontenc}
\usepackage{hyperref}
\title{Zeeuws Language Resources for K-12 Curriculum Development}
\author{Compiled by Claude Sonnet 5.5 (claude-sonnet-5-5) for Andr\'as Kornai}
\date{}
\begin{document}
\maketitle

\section{Introduction}
Zeeuws (Zeelandic, ISO 639-3 zea) is spoken in the Dutch province of Zeeland, with an estimated 220,000 speakers. It is usually classed as a Low Franconian variety close to Dutch. The Dutch government has declined to recognise it as a regional language under the European Charter. The searches found no evidence of Zeeuws being taught in schools and no school materials. This report lists only what the searches turned up.

\section{Primary Educational Resources}
No primers, textbooks or school programmes for Zeeuws turned up in the searches.

\section{Linguistic and Reference Materials}
\begin{itemize}
\item Zeelandic, Wikipedia overview of the language, its status and its spelling: \url{https://en.wikipedia.org/wiki/Zeelandic}
\item Omniglot, Zeelandic (Ze\^euws): notes on writing and mentions a Zeelandic-Dutch dictionary, without giving its details: \url{https://www.omniglot.com/writing/zeelandic.htm}
\item \emph{Woordenboek der Zeeuwse dialecten}, the standard dictionary of the Zeeland dialects, compiled by the Zeeuwse Vereniging voor Dialectonderzoek under Dr. Ghijsen and completed in 1964. A review of its 2003 supplement is at: \url{https://dbnl.org/tekst/_bie001200301_01/_bie001200301_01_0013.php}
\item Taal en Tongval (1979) article discussing Zeeland dialect notation and the Ghijsen dictionary: \url{https://www.dbnl.org/tekst/_taa007197901_01/_taa007197901_01_0024.php}
\end{itemize}

\section{Religious and Community Resources}
No religious or community-organisation resources were found in the searches.

\section{Digital and Audio Resources}
No audio, video or software resources were found. The Lexvo language page is only a catalogue entry: \url{https://finto.fi/lexvo/en/page/zea}

\section{Curriculum Development Context}
The sources reviewed describe Zeeuws as differing from standard Dutch mainly in pronunciation, but also in grammar and vocabulary. Official recognition has been refused. A curriculum would therefore have to build its own materials, using the dialect dictionary and its supplement as the lexical base. Current school or community initiatives could not be confirmed from the searches. The Zeeuwse Vereniging voor Dialectonderzoek would be a sensible organisation to contact.
\end{document}